\section*{अध्याय \ref{ch:b3:bioinformatics} --- जैवसूचना-विज्ञान और अनुक्रम विश्लेषण}

\begin{solution}{exo:b3:bioinformatics:1}
वैश्विक: दोनों अनुक्रम सिरे से सिरे तक संरेखित, हर अवशेष किसी स्तंभ में
--- ऐसे दो प्रोटीनों के लिए जो अपनी पूरी लंबाई पर समजात माने जाते हों,
जैसे किसी गृहव्यवस्था एंज़ाइम के \omterm{def:b3:genomics:comparative}{ऑर्थोलॉग}। स्थानिक: उपअनुक्रमों का
सर्वोत्तम-अंक युग्म, शेष अनदेखा --- किसी साझा प्रांत को ढूँढ़ने के लिए
(दो अन्यथा असंबंधित संकेतन प्रोटीनों में कोई SH2 प्रांत), या किसी लंबे
जीनोमी अनुक्रम में कोई जीन ढूँढ़ने के लिए।
\end{solution}

\begin{solution}{exo:b3:bioinformatics:2}
दोनों ओर सीमाएँ $0,-1,-2,-3$। पंक्ति A: $1, 0, -1$। पंक्ति G: $0, 0, -1$।
पंक्ति C: $-1, -1, 1$। सर्वोत्तम $F(3,3) = 1$: AAC के ऊपर AGC, बिना
किसी रिक्ति (मेल, असंगति, मेल: $1 - 1 + 1 = 1$)।
\end{solution}

\begin{solution}{exo:b3:bioinformatics:3}
$s(a,a) = \lambda^{-1}\log\bigl(q_{aa}/p_{a}^{2}\bigr)$। ट्रिप्टोफ़ैन
दुर्लभ है ($p_{W} \approx 0.013$), इसलिए दो ट्रिप्टोफ़ैनों के यादृच्छिक
रूप से संरेखित होने की संभावना, $p_{W}^{2}$, बहुत कम है, और कोई
संरक्षित ट्रिप्टोफ़ैन युग्म समजातता का किसी संरक्षित ल्यूसीन युग्म ($p_{L}
\approx 0.1$) से कहीं प्रबल चिह्न है; लघुगणकीय-संभावना अनुपात उसी के
अनुसार बड़ा है।
\end{solution}

\begin{solution}{exo:b3:bioinformatics:4}
E-मान उन \omterm{def:b3:bioinformatics:alignment}{संरेखणों} की वह संख्या है जिनका अंक कम-से-कम उतना ही ऊँचा हो
और जो इसी पृच्छा की इसी आकार के आँकड़ाकोश के सामने की गई खोज में संयोग
से प्रत्याशित हों। $E = 3$ का अर्थ है कि संयोग से ऐसे तीन अंक प्रत्याशित
हैं: वह मेल समजातता का प्रमाण नहीं है (हो सकता है वह समजात हो ही, पर
खोज यह बता नहीं सकती)।
\end{solution}

\begin{solution}{exo:b3:bioinformatics:5}
$mn = 400\times 2\times 10^{11} = 8\times 10^{13} = 2^{46.2}$। $E(45) =
2^{1.2} \approx 2.3$; $E(55) = 2^{-8.8} \approx 2\times 10^{-3}$;
$E(65) = 2^{-18.8} \approx 2\times 10^{-6}$। $E = 10^{-3}$ के लिए $S' =
46.2 + 10.0 = 56$ बिट चाहिए। \num{40}-अवशेष पृच्छा का $mn$ दस गुना
छोटा, $2^{42.9}$ है: $53$ बिट पर्याप्त हैं --- पर कोई छोटी पृच्छा उतने
तक भी शायद ही पहुँचती है।
\end{solution}

\begin{solution}{exo:b3:bioinformatics:6}
\omterm{def:b3:bioinformatics:motif}{सूचना-अंश}: $2$, $1$, $0$, और $2 - H$, जिसमें $H = -(0.7\log_{2}
0.7 + 3\times 0.1\log_{2} 0.1) = 0.36 + 1.00 = 1.36$ है, इसलिए $0.64$। कुल
$R = 3.64$ बिट। संयोगी मेल: दो लड़ियों पर $9.2\times 10^{6}$ स्थितियाँ
$\times 2^{-3.64} \approx 7\times 10^{5}$ --- \omterm{def:b3:bioinformatics:motif}{अभिप्ररूप} अकेला लगभग बेकार है।
\end{solution}

\begin{solution}{exo:b3:bioinformatics:7}
कई अवशेषों का कोई अंतर्वेशन एक ही उत्परिवर्तनी घटना है, इसलिए उसकी लागत
उसकी लंबाई के साथ रैखिक रूप से नहीं बढ़नी चाहिए: कोई खोलने की लागत और
कोई छोटी बढ़ाने की लागत इसका प्रतिरूप बनाती हैं। बहुत ऊँचा खोलने का दंड
वहाँ असंगतियाँ थोप देता है जहाँ कोई रिक्ति चाहिए और किसी सच्चे
अंतर्वेशन के बाद सब कुछ गलत संरेखित कर देता है; बहुत नीचा दंड हर जगह
रिक्तियाँ बिखेर देता है, अवशेषों को संयोग से मिला देता है और समरूपता
फुला देता है।
\end{solution}

\begin{solution}{exo:b3:bioinformatics:8}
ज़रूरी नहीं: \omterm{def:b3:bioinformatics:alignment}{संरेखण} \num{130}-अवशेष के किसी खंड को ढकता है, जो किसी
साझा प्रांत का लक्षण है, न कि अपनी लंबाई पर संरेखित किसी \omterm{def:b3:genomics:comparative}{ऑर्थोलॉग} का।
जाँच: मक्खी प्रोटीन को वापस मानव प्रोटीनांश के सामने खोजिए (क्या पृच्छा
उसका सर्वोत्तम मेल है, पूरी लंबाई पर?), प्रांत को किसी \omterm{def:b3:bioinformatics:msa}{प्रोफ़ाइल} HMM से
पहचानिए, और कई जातियों में उस कुल का जीन-वृक्ष बनाइए।
\end{solution}

\begin{solution}{exo:b3:bioinformatics:9}
कोई \omterm{def:b3:bioinformatics:msa}{प्रोफ़ाइल} स्तंभ-दर-स्तंभ अंकित करती है कि कुल क्या सहता है: ऐसी
स्थिति जो सदा जलविरोधी रहती है पर कभी एक ही अवशेष नहीं, कोई अचर
उत्प्रेरक अवशेष, ऐसी स्थिति जो आधे कुल में सदा रिक्ति रहती है। कोई
युग्मवार \omterm{def:b3:bioinformatics:alignment}{संरेखण} हर अवशेष को किसी एक दूसरे अवशेष के सामने अंक देता है और
यह जान नहीं सकता कि स्थिति 40 का वैलीन पृच्छा के वहाँ पड़े आइसोल्यूसीन
जितना ही अच्छा है। \omterm{def:b3:bioinformatics:msa}{प्रोफ़ाइल} संरक्षित स्तंभों को भार भी देती है, इसलिए
जहाँ कुल संरक्षित है वहीं संकेंद्रित कोई कमज़ोर सादृश्य सार्थक बन जाता है।
\end{solution}

\begin{solution}{exo:b3:bioinformatics:10}
प्रति स्तंभ धनात्मक \omterm{prop:b3:bioinformatics:logodds}{प्रत्याशित अंक} $\mu > 0$ के साथ, दो यादृच्छिक
अनुक्रमों के विकर्ण पर संचयी अंक धनात्मक अपवाह वाला यादृच्छिक चरण-पथ
है: $n$ स्तंभों के बाद वह लगभग $\mu n$ होता है, इसलिए सर्वोत्तम
\omterm{def:b3:bioinformatics:alignment}{स्थानिक संरेखण} मूलतः पूरा ही होता है और उसका अंक $\mu n$ की तरह बढ़ता
है, $\log n$ की तरह नहीं। कार्लिन--ऑल्टशुल सिद्धांत, जिसे ऋणात्मक अपवाह
चाहिए ताकि ऊँचे अंक दुर्लभ भ्रमण हों, लागू नहीं होता और कोई $\lambda$
नहीं होता। \qty{60}{\%} GC पर DNA के लिए मेल की
संभावना $2(0.3^{2}) + 2(0.2^{2}) = 0.26$ है, इसलिए \omterm{prop:b3:bioinformatics:logodds}{प्रत्याशित अंक}
$0.26 - 0.74 = -0.48$ है: अब भी ऋणात्मक, और सांख्यिकी चलती
है; पर मेल $+1$, असंगति $-0.3$ जैसी किसी योजना की प्रत्याशा
$+0.04$ होती और वह पूरे जीनोम को एक ही \omterm{def:b3:bioinformatics:alignment}{संरेखण} बता देती।
\end{solution}

\begin{solution}{exo:b3:bioinformatics:11}
\omterm{def:b3:bioinformatics:blast}{बीज} और शृंखलन: किसी हैश सारणी से दोनों अनुक्रमों के बीच $k$-मरों के
ठीक-ठीक या लगभग ठीक मेल ढूँढ़िए, उन्हीं को रखिए जो संगत विकर्णों पर
पंक्तिबद्ध हों, उन्हें शृंखलित कीजिए, और \omterm{thm:b3:bioinformatics:nw}{गतिक क्रमादेशन} केवल शृंखलित
बीजों के बीच के अंतरालों में चलाइए; इसमें उन क्षेत्रों के \omterm{def:b3:bioinformatics:alignment}{संरेखण} छूट
जाते हैं जहाँ कोई \omterm{def:b3:bioinformatics:blast}{बीज} न हो (बहुत अपसरित खंड)। पट्टीकरण: यदि दोनों
अनुक्रम लगभग सहरेखी माने जाते हों, तो केवल विकर्ण के चारों ओर चौड़ाई
$w$ की पट्टी के भीतर के कोष्ठक निकालिए, $wn$ लागत पर, $mn$ के बजाय;
इसमें पट्टी से बड़े किसी भी अंतर्वेशन वाला \omterm{def:b3:bioinformatics:alignment}{संरेखण} छूट जाता है।
\end{solution}

\begin{solution}{exo:b3:bioinformatics:12}
कूटलेखी अनुक्रम का आवर्त तीन का होता है: तीनों कूटक-स्थितियों के
क्षार-संघटन भिन्न होते हैं (तीसरी सबसे अभिनत होती है), और हर जाति में
कूटक-उपयोग असमान होता है। क्रम में तीन कूटलेखी अवस्थाओं वाला कोई
प्रतिरूप, जिनमें से हर एक उसी कूटक-स्थिति के संघटन से क्षार उत्सर्जित
करती है, कुछ दर्जन कूटकों की किसी खिड़की पर कूटलेखी DNA को अकूटलेखी
अवस्था की तुलना में अधिक प्रायिकता देता है, विराम-कूटकों के बिना भी।
मानव जीनोम में बहिर्वेशन छोटे होते हैं (\qty{150}{bp}), जो किलोक्षारों
के अंतर्वेशनों से अलग होते हैं, इसलिए कूटलेखी संकेत छोटा और बीच-बीच में
टूटा होता है; प्रतिरूप को विच्छेदन स्थल भी पहचानने पड़ते हैं, जो कमज़ोर
संकेत हैं, और अकूटलेखी अनुक्रम की भारी मात्रा अनेक मिथ्या कूटलेखी खंड
पैदा करती है।
\end{solution}

\begin{solution}{pb:b3:bioinformatics:1}
\textbf{1.} पंक्तियाँ K, Q, T; स्तंभ K, A, Q, T; सीमाएँ $0,-1,-2,-3,-4$
और $0,-1,-2,-3$। पंक्ति K: $1, 0, -1, -2$; पंक्ति Q: $0, 0, 1, 0$;
पंक्ति T: $-1, -1, 0, 2$। सर्वोत्तम $2$: \texttt{KAQT} के ऊपर \texttt{K-QT}।
\textbf{2.} सर्वोत्तम स्थानिक अंक $3$: \texttt{CAT} के सामने \texttt{CAT}
(GATCAT के अवशेष 4--6 और ACAT के 2--4); \texttt{A-CAT} के सामने
\texttt{ATCAT} का अंक भी $4 - 1 = 3$ है।
\textbf{3.} $300\times 450 = 1.35\times 10^{5}$ अद्यतन; आँकड़ाकोश के
सामने $300\times 1.2\times 10^{11} = 3.6\times 10^{13}$।
\textbf{4.} $3.6\times 10^{4}$ s, प्रति पृच्छा दस घंटे; \omterm{def:b3:bioinformatics:blast}{BLAST} के \omterm{def:b3:bioinformatics:blast}{बीज}
सारणी का लगभग पूरा भाग छोड़ देते हैं और सेकंडों में उत्तर देते हैं।
\textbf{5.} $q_{ab} = p_{a}p_{b}e^{\lambda s}$। ऐलानीन: $e^{1.39} = 4.0$,
$q_{AA} = 0.074^{2}\times 4.0 = 0.022$, अनुपात $4$। ट्रिप्टोफ़ैन:
$e^{3.82} = 45$, $q_{WW} = 0.013^{2}\times 45 = 0.0077$, अनुपात $45$।
कोई संरेखित ट्रिप्टोफ़ैन युग्म समजातों में संयोग से $45$ गुना अधिक
बारंबार है, कोई ऐलानीन युग्म केवल चार गुना; फिर भी परम रूप में ऐलानीन
युग्म अधिक आम हैं क्योंकि ऐलानीन आम है।
\textbf{6.} \qty{24}{\%} गोधूलि क्षेत्र में पड़ता है, जहाँ यादृच्छिक
\omterm{def:b3:bioinformatics:alignment}{संरेखण} \qtyrange{15}{20}{\%} तक पहुँच जाते हैं; \omterm{def:b3:bioinformatics:alignment}{संरेखण} का E-मान, सही
स्थितियों पर संरक्षित \omterm{def:b3:bioinformatics:motif}{अभिप्ररूप}, किसी ज्ञात कुल से प्रोफ़ाइल-HMM मेल,
या कोई साझा वलन इसका निपटारा कर देगा।
\textbf{7.} $mn = 300\times 1.2\times 10^{11} = 3.6\times 10^{13}$;
$\log_{2}(mn) = 45.0$।
\textbf{8.} $E(92) = 2^{45 - 92} = 2^{-47} \approx 7\times 10^{-15}$;
$E(38) = 2^{7} = 128$।
\textbf{9.} $E = 10^{-3}$ पर $S' = 45 + 10 = 55$ बिट; $E = 1$ पर $45$
बिट।
\textbf{10.} $mn$ दस गुना: $E \approx 7\times 10^{-14}$, फिर भी
भारी-भरकम।
\textbf{11.} $E = 128$ के साथ दसवाँ मेल वही है जो संयोग पैदा करता है;
$60$ अवशेषों पर \qty{40}{\%} समरूपता कोई प्रमाण नहीं है। अवशेष 200--260
की प्रांत आँकड़ाकोश के सामने कोई प्रोफ़ाइल-खोज दिखा सकती है कि वह खंड
कोई ज्ञात प्रांत है या नहीं, ऐसी सांख्यिकी के साथ जो युग्मवार तुलना के
पास नहीं।
\textbf{12.} पूरी लंबाई पर पारस्परिक सर्वोत्तम मेल एक-से-एक ऑर्थोलॉगी
के अनुकूल हैं; वे उसे सिद्ध नहीं करते --- विभाजन के बाद किसी एक
वंशक्रम में हुआ कोई द्विगुणन दो \omterm{def:b3:genomics:comparative}{सह-ऑर्थोलॉग} देता है, और सच्चे
\omterm{def:b3:genomics:comparative}{ऑर्थोलॉग} के खो जाने पर कोई \omterm{def:b3:genomics:comparative}{पैरालॉग} सर्वोत्तम मेल रह सकता है। कई
जातियों वाला कोई जीन-वृक्ष ही कसौटी है।
\textbf{13.} $R = 2 + 2 + 1.6 + 2 + 0.8 + 1.2 + 0.4 + 0.3 = 10.3$ बिट।
\textbf{14.} $3.2\times 10^{9}\times 2^{-10.3} \approx 2.5\times 10^{6}$
संयोगी मेल।
\textbf{15.} $\log_{2}(3.2\times 10^{9}/200) = \log_{2}(1.6\times 10^{7})
\approx 24$ बिट।
\textbf{16.} नियत अंतराल पर साथ आना बिट जोड़ देता है: $10.3 + 8 =
18.3$, जो $8$ पूरे करता है, छूटे हुए $13.7$ में से; लगभग $5.7$ बिट
(संयोगी मेलों में $50$ का गुणक) कहीं और से आने चाहिए --- \omterm{def:b3:chromatin-epigenetics:nucleosome}{क्रोमैटिन} की
सुलभता, और साझेदार।
\textbf{17.} $H = -(0.5\log_{2}0.5 + 0.5\log_{2}0.5) = 1$ बिट; $R = 2 -
1 = 1$ बिट।
\textbf{18.} किसी जीवाणु कारक को अपने थोड़े-से स्थल किसी
\qty{4.6}{Mb} जीनोम में \omterm{def:b3:chromatin-epigenetics:nucleosome}{क्रोमैटिन} की किसी मदद के बिना ढूँढ़ने पड़ते हैं:
उसे कोई $19$ बिट चाहिए, और उसके स्थल लंबे तथा संरक्षित होते हैं। कोई
सुकेंद्रकी जीनोम हज़ार गुना बड़ा है और उसे दस बिट और चाहिए, फिर भी उसके
कारकों के स्थल छोटे हैं; वे विशिष्टता संयोजन से और सुलभ \omterm{def:b3:chromatin-epigenetics:nucleosome}{क्रोमैटिन} की
सीमा से पाते हैं, जो नियमन को अधिक विकासनीय भी बनाता है, क्योंकि कोई
छोटा स्थल सहज ही पाया या खोया जाता है।
\textbf{19.} $3/64 = 0.047$; किसी विराम-कूटक से पहले कूटकों की
प्रत्याशित संख्या $64/3 \approx 21$ है।
\textbf{20.} $p_{A} = p_{T} = 0.31$, $p_{G} = p_{C} = 0.19$: $P(\text{TAA})
= 0.31^{3} = 0.030$, $P(\text{TAG}) = P(\text{TGA}) = 0.31^{2}\times 0.19
= 0.018$; कुल $0.066$, प्रत्याशित ढाँचा-लंबाई $15$ कूटक।
AT-समृद्ध DNA विराम-कूटकों से भरा है, इसलिए यादृच्छिक खुले ढाँचे छोटे
होते हैं और लंबे ढाँचे अधिक अलग दिखते हैं।
\textbf{21.} $(61/64)^{100} = e^{-4.80} = 0.008$; $(61/64)^{300} =
e^{-14.4} = 5.6\times 10^{-7}$।
\textbf{22.} हर उच्चिष्ठ खुला ढाँचा किसी विराम-कूटक पर समाप्त होता है,
और $3.2\times
10^{9}$ कूटक-आरंभों में $3.2\times 10^{9}\times 3/64 = 1.5\times
10^{8}$ विराम-कूटक हैं: लगभग $1.5\times 10^{8}\times 0.008 = 1.2\times 10^{6}$
यादृच्छिक ढाँचे कम-से-कम $100$ कूटकों के, और $1.5\times 10^{8}\times
5.6\times 10^{-7} \approx 80$ कम-से-कम $300$ के।
\textbf{23.} \qty{4.6}{Mb} के किसी जीवाणु में कोई $4\times 10^{5}$
विराम-कूटक हैं और इसलिए $100$ कूटकों के लगभग \num{3500} संयोगी ढाँचे,
पर $300$ के लगभग कोई नहीं; उसके जीन औसतन $300$ कूटक के हैं और DNA का
\qty{88}{\%} कूटलेखी है, इसलिए कोई लंबा खुला ढाँचा लगभग सदा कोई जीन
होता है। कृमि में DNA का \qty{1.5}{\%} कूटलेखन करता है, बहिर्वेशन
औसतन $50$ कूटक के हैं --- संयोग की देहली से भी छोटे --- और $100$
कूटकों के दस लाख यादृच्छिक ढाँचे उन्हें दबा देते हैं। सुकेंद्रकी
जीन-खोजी विच्छेदन-स्थल संकेत, किसी गुप्त मार्कोव प्रतिरूप में
कूटक-अभिनति, ज्ञात प्रोटीनों से समजातता और, सबसे बढ़कर, अनुक्रमित
अनुलेख काम में लेते हैं।
\textbf{24.} किसी विच्छेदित संदेशवाहक से आया \omterm{def:b3:genomics:ngs}{वाचन} जीनोम से दो टुकड़ों
में संरेखित होता है, जिनके बीच कोई अंतर्वेशन है: वह विभाजन दोनों
विच्छेदन स्थलों को क्षार तक चिह्नित कर देता है; वाचन-आच्छादन
बहिर्वेशनों की सीमा खींचता है और युग्मित \omterm{def:b3:genomics:ngs}{वाचन} क्रमिक बहिर्वेशनों को
एक ही अनुलेख में जोड़ देते हैं, जिससे तय हो जाता है कि कई संभावित
विच्छेदन स्थलों में से कौन-सा काम में आया।
\textbf{25.} सर्वोत्तम मेल के लिए $E \approx 7\times 10^{-15}$; लगभग
$24$ बिट, जीनोम में $200$ स्थल निर्दिष्ट करने को; और पूरे जीनोम में कोई
$80$ संयोगी वाचन-ढाँचे, $300$ कूटकों के।
\end{solution}
